% Fragment: regularity, branching, fusion-ring reduction, and characters.
% The main document supplies theorem environments, reviewnote, and artifact.

\section{Regularity from a positive conformal-vector identity}
\label{sec:regularity}

This section concerns the actual PCT-stable subalgebra $U\subset B$ defined
in the construction. Its conclusions do not require properness. We use
the convention $Y(a,z)=\sum_{n\in\mathbb Z}a_nz^{-n-1}$ and write
$L^e(n)=e_{n+1}$ for the Virasoro operators of a conformal vector $e$.
The input from the local construction is the identity
\begin{equation}
 \omega_U=\omega_{A_{1,1}}
   +\frac{3403}{4563}\omega_{A_{1,39}}
   +\frac{25}{6}\overline p+\frac{25}{26}\overline z.
 \label{eq:positive-stress}
\end{equation}
Here $p$ and $z$ are the conformal vectors of actual embedded copies of
the three-state Potts algebra $K(\mathfrak{sl}_2,3)$ and the parafermion
algebra $K(\mathfrak{sl}_2,13)$; a bar denotes normalized averaging over
the compact current group. Each subalgebra is regular and unitary, and
its intrinsic grading agrees with the ambient grading on its states.
The averages take place in the finite-dimensional space $U_2$.
Finite-dimensional convexity replaces each average by a finite convex
combination of conjugate conformal vectors, so that
\begin{equation}
 \omega_U=\sum_{i=1}^N c_i e_i,\qquad c_i>0.
 \label{eq:finite-positive-stress}
\end{equation}
Zero coefficients, if any, are omitted. All simple modules of each
regular subalgebra $A_i$ have nonnegative lowest weights, with zero
occurring only for its vacuum module.

\begin{reviewnote}{R1: the embedded regular factors are local inputs}
The existence and normalization of the Potts and parafermion fields,
their inclusion in the generated $U$, and
\eqref{eq:positive-stress} are imported from the earlier construction
and the proof records
\artifact{hi_cft_20260923_next_structure_c2_energy_bound.md} and
\artifact{hi_cft_20260923_next_structure_top_bound.md}.
They are not consequences of the finite closure theorem alone.
This manuscript gives the deductions from those inputs, but does not
reproduce their full field-identification calculation. A review of
regularity must verify these embeddings, not only the displayed
positive coefficients.
\end{reviewnote}

\begin{lemma}[Orthogonal energy bound]
Let $V$ be a unitary VOA of CFT type admitting an identity of the form
\eqref{eq:finite-positive-stress}, with the properties just stated.
Then $V$ is $C_2$-cofinite.
\end{lemma}
\begin{proof}
For an ordinary simple $A_i$-module $M$, put
\[
 C_{2,A_i}(M)=\operatorname{span}\{a_{-2}m:a\in A_i,\ m\in M\}.
\]
These quotients are finite-dimensional by
\href{https://arxiv.org/html/math/0204021#S5}{Abe--Buhl--Dong,
Proposition~5.2}. There are finitely many simple $A_i$-modules. Hence
there is a finite number $D_i$ bounding the intrinsic energy in all
their $C_{2,A_i}$ quotients; this includes their lowest weights.
Regularity decomposes the weak $A_i$-module $V$ into ordinary simples,
and the unitary form permits an orthogonal decomposition. Since
\[
 C_{2,A_i}(V)\subset C_2(V),
\]
the self-adjoint intrinsic energy satisfies
\[
 \langle L^{e_i}(0)v,v\rangle\le D_i\|v\|^2
 \qquad (v\in C_2(V)^\perp).
\]
For $v\in V_n\cap C_2(V)^\perp$, the conformal-vector identity gives
\[
 n\|v\|^2
 =\sum_i c_i\langle L^{e_i}(0)v,v\rangle
 \le\Bigl(\sum_i c_iD_i\Bigr)\|v\|^2.
\]
Thus the graded quotient $V/C_2(V)$ vanishes above a finite weight.
No simultaneous diagonalization of the different energy operators is
required. The same proof applies directly to averages, since every
conjugate factor has the same bound $D_i$.
\end{proof}

\begin{lemma}[Lowest weights and projectivity of the vacuum]
Every simple ordinary $U$-module has nonnegative lowest conformal weight,
and the only simple module of lowest weight zero is $U$. Moreover $U$
is projective in the category of grading-restricted generalized modules.
\end{lemma}
\begin{proof}
Let $E$ be the lowest space of a simple ordinary module of weight $h$.
Every mode of an $A_i$ state that lowers intrinsic weight also lowers
ambient weight, and therefore annihilates $E$. Thus
$E\subset\Omega_{A_i}(M)$, the Zhu highest space. Restriction initially
gives a weak $A_i$-module; regularity decomposes it into ordinary simple
modules. The highest space of an ordinary simple is its top, so each
$L^{e_i}(0)|_E$ is diagonalizable with nonnegative real spectrum.
Taking traces, rather than assuming the operators commute, gives
\begin{equation}
 h\dim E=\sum_i c_i\operatorname{Tr}_E L^{e_i}(0)\ge0.
 \label{eq:top-trace}
\end{equation}
If $h=0$, every summand vanishes. Hence every partial energy is zero
on $E$, which lies in sums of vacuum tops for every $A_i$. Consequently
$L^{e_i}(-1)E=0$, and \eqref{eq:finite-positive-stress} gives
$L_U(-1)E=0$. Any nonzero $v\in E$ is vacuum-like: translation covariance
and lower truncation show that $Y_M(a,z)v$ has no negative powers of $z$.
The map $a\mapsto a_{-1}v$ is a nonzero module map $U\to M$, hence an
isomorphism.

Now let $\pi:E\twoheadrightarrow U$ be an epimorphism in the generalized
module category. $C_2$-cofiniteness gives finite length, and the preceding
argument excludes negative weights in all its composition factors.
The finite generalized weight-zero space $E_{[0]}$ therefore has no
lower-weight target for a lowering mode. It is again in every
$\Omega_{A_i}(E)$. On this space $L_U(0)$ is nilpotent, so
\[
 0=\operatorname{Tr}_{E_{[0]}}L_U(0)
  =\sum_i c_i\operatorname{Tr}_{E_{[0]}}L^{e_i}(0).
\]
Diagonalizability and nonnegativity of each partial energy imply it is
zero. The vacuum-top argument gives $L_U(-1)E_{[0]}=0$. Choose
$v\in E_{[0]}$ with $\pi(v)=\mathbf1$. Then
\[
 s:U\longrightarrow E,\qquad s(a)=a_{-1}v,
 \qquad \pi s=\mathrm{id}_U,
\]
is a section. Every epimorphism onto the vacuum splits.
\end{proof}

\begin{proposition}[Regularity of the actual generated algebra]
The VOA $U$ is strongly rational.
\end{proposition}
\begin{proof}
Simplicity, CFT type, and self-contragredience follow from the unitary
construction. The preceding lemmas supply $C_2$-cofiniteness and a
projective vacuum. The category $\mathcal C$ of grading-restricted
generalized modules is finite abelian and braided, with right-exact
tensor product and Grothendieck--Verdier duality whose dualizing object
is $U'=U$; the relevant VOA tensor-category hypotheses are reviewed in
\href{https://arxiv.org/html/2108.01898}{McRae, Sections~2.3 and~3}.

Here is the remaining categorical argument. Realize $\mathcal C$ as
finite-dimensional modules for a finite-dimensional algebra, denote
underlying vector-space dimension by $d_0$, and let $P$ correspond to
its regular module. Every $Z$ is a quotient of $P^{\oplus d_0(Z)}$.
For fixed $X$, right exactness gives
\[
 d_0(Z\otimes X)\le C_Xd_0(Z),\qquad
 \dim\operatorname{End}(X^{\otimes n})
 \le d_0(\mathbf1)^2C_X^{2n}<n!
\]
for large $n$, where $C_X=\max\{1,d_0(P\otimes X)\}$.
For nonzero indecomposable $X$, Grothendieck--Verdier evaluation
$DX\otimes X\to\mathbf1$ is nonzero, hence surjective onto the simple
unit, and splits by projectivity. Thus $X$ is non-negligible and has
moderate growth in the precise sense of
\href{https://arxiv.org/html/2412.17681v3}{Etingof--Penneys,
Theorem~1.1}; that theorem proves $X$ rigid. Finite direct sums give
rigidity of all objects. Tensoring is now exact, and
\[
 \operatorname{Hom}(X,-)
 \cong\operatorname{Hom}(\mathbf1,-\otimes X^*)
\]
is exact. Every object is projective, so $\mathcal C$ is semisimple and
$U$ is rational. This does not assume complete unitarity of its modules.
The local argument is recorded in
\artifact{hi_cft_20260923_next_structure_rationality.md}.
\end{proof}

\section{Branching and the boundary fusion ring}
\label{sec:branching}

From this point we use the properness conclusion and $\dim U_1=6$ of
the finite closure argument. Before that conclusion is supplied, every
statement in this section is conditional on the proper branch.

\subsection{A universal top bound and actual lattice branching}

The sharper bound uses an actual regular subalgebra $A_L\subset U$.
In the notation of the local construction, a rational lattice--Potts
core $\mathcal T$ and $K(\mathfrak{sl}_2,13)$ have pointed charges
$(j,k)\in(\mathbb Z/13)^2$ with quadratic form
$-(j^2+k^2)/13$. Exact weight-two vectors witness the simple-current
charge $(5,1)$ (up to reversing the second charge). Locality and
multiplicity one in a connected extension force its entire isotropic
line, yielding
\[
 A_L=\bigoplus_{k=-5j}\mathcal T_{1,j}\otimes\mathrm{PF}_{0,k}
 \subset U,\qquad
 \operatorname{Rep}(A_L)\simeq
 \mathrm{Fib}\boxtimes\mathrm{SO}(3)_{13}.
\]
The field witness and extension argument are in
\artifact{hi_cft_20260923_next_structure_order13_core.md} and
\artifact{hi_cft_20260923_next_structure_core_extension.json}.
There are fourteen simple tops, labelled $x_a=(1,a)$ and
$y_a=(\tau,a)$ for $0\le a\le6$.

Let $E$ be any simple $U$-module top and write $x_a,y_a$ also for the
nonnegative integer multiplicities of these $A_L$ tops in $E$.
Set $D=\dim E$. For each top, let $L$ be the trace of the left affine
energy, $Q=\operatorname{Tr}H_R(0)^2$, and $P,Z$ the Potts and
parafermion energy traces. Its contribution to \eqref{eq:top-trace} is
\begin{equation}
 W=L+\frac{83}{6084}Q+\frac{25}{6}P+\frac{25}{26}Z.
 \label{eq:top-W}
\end{equation}
This is a trace formula on a complete right $\mathfrak{sl}_2$
representation, not an assertion about a single partial-energy
eigenvector. The exact dimensions and trace contributions are
\[
\begin{array}{c|r|r@{\qquad}c|r|r}
 \text{top}&d&4680W&\text{top}&d&4680W\\\hline
 x_0&1&0&y_0&1&7800\\
 x_1&1&600&y_1&7&36960\\
 x_2&5&10440&y_2&15&110880\\
 x_3&9&56640&y_3&2&4080\\
 x_4&2&5520&y_4&14&81480\\
 x_5&14&148560&y_5&6&18240\\
 x_6&6&54240&y_6&2&4920
\end{array}
\]
Complete right $\mathfrak{sl}_2$ characters imply
$y_2\ge x_5$ and $y_4\ge x_5+x_6$: these compare charges $10,12$
and charges $7,9$, respectively, with the common left doublet removed
in the second comparison. Direct comparison of the displayed table gives
\begin{align}
 4680\bigl(\tfrac{17}{10}D-W_{\rm total}\bigr)
 ={}&8460(y_2-x_5)+28716(y_4-x_5-x_6)\notag\\
 &+7956x_0+7356x_1+29340x_2+14964x_3+10392x_4\notag\\
 &+22212x_6+156y_0+18732y_1+11832y_3\notag\\
 &+1188y_4+29496y_5+10992y_6\ge0.
 \label{eq:sharp-top-bound}
\end{align}
Since $W_{\rm total}=hD$, this proves $0\le h\le17/10<2$ for
\emph{every} simple $U$-module, without assuming it is unitarizable.

\begin{reviewnote}{R2: the core-top table and finite-top actions}
The inequality \eqref{eq:sharp-top-bound} can be checked directly from
the displayed table. Deriving that table, including all minimizing
charge representatives, uses
\artifact{hi_cft_20260923_next_structure_order13_top_bound.md}.
The construction of $\mathcal T$, the thirteen-charge extension, and
completeness of the resulting fourteen-module list are substantial
local inputs whose full derivations are not repeated here. Likewise,
the finite-top irreducibility used below is imported from exact field
actions, particularly
\artifact{hi_cft_20260923_next_81top_all.json}; it should not be replaced
by numerical cyclicity alone.
\end{reviewnote}

The strongly rational restriction $B\downarrow_U$ is semisimple. Its
simple summands have integer lowest weights, hence only weights $0,1$
are possible. The vacuum occurs once, while the complementary current
space has dimension $86-6=80$. Its certified irreducible zero-mode
action therefore implies
\begin{equation}
 B\downarrow_U=U\oplus X,\qquad h_X=1,\qquad\dim X_{\rm top}=80.
 \label{eq:B-branching}
\end{equation}
The four mixed nontrivial $E_6\oplus A_2$ lattice sectors start at
weight $1$, with $81$-dimensional tops. Their next possible top weight
would be $2$, excluded by \eqref{eq:sharp-top-bound}. Their exact
top-action certificates use a one-dimensional Cartan eigenline,
positive-metric adjoint identities, and a cyclic span of dimension
$81$. The adjoint identities make the generated operator algebra a
star algebra; orthogonal complements are therefore invariant, and
the unique eigenline plus cyclicity proves irreducibility.

Write the actual lattice extension as $H=B\oplus J\oplus J^2$. Its
restriction is consequently
\begin{equation}
 H\downarrow_U=\mathbf1\oplus X\oplus Y\oplus Y^*,\qquad
 h_Y=1,\quad\dim Y_{\rm top}=81.
 \label{eq:H-branching-fork}
\end{equation}
At this stage $Y\cong Y^*$ has not been assumed. The three types
$\mathbf1,X,Y$ are distinguished by their lowest spaces.

\subsection{The rank-four or rank-six alternative}

Put $\mathcal C=\operatorname{Rep}(U)$ and
$\mathcal F=\mathcal C_H$, the category of \emph{all} right modules
over the extension algebra $H$. The strict positivity of nonvacuum
lowest weights gives the positive spherical structure through character
asymptotics. Since $H$ is holomorphic, it is Lagrangian and the
\href{https://arxiv.org/abs/1009.2117}{DMNO center correspondence}
gives $\mathcal C\simeq Z(\mathcal F)$.
Write $s=\operatorname{FPdim}_{\mathcal C}B$. Extension dimension
identities give
\[
 d_X=s-1,\qquad d_Y=d_{Y^*}=s,\qquad
 \operatorname{Dim}\mathcal C=9s^2,\quad
 \operatorname{Dim}\mathcal F=3s.
\]
Pseudo-unitarity of $\mathcal F$ follows from that of its center and
the squared global-dimension and Frobenius--Perron-dimension identities.
No $C^*$ structure is needed in this section.

For center induction $I$ and forgetful functor $F$, one has $H=I(\mathbf1)$
and $FI(\mathbf1)=\bigoplus_{a\in\operatorname{Irr}\mathcal F}a\otimes a^*$.
Adjunction therefore gives
\begin{equation}
 \operatorname{rank}\mathcal F
 =\dim\operatorname{End}_{\mathcal C}(H)
 =\begin{cases}4,&Y\not\cong Y^*,\\6,&Y\cong Y^*.
 \end{cases}
 \label{eq:rank-fork}
\end{equation}
The actual order-three automorphisms of $H/B$ inject into
$\operatorname{Inv}(\mathcal F)$. This injection is already spherical:
under the composition/convolution description of
$\operatorname{End}(I(\mathbf1))$ in
\href{https://arxiv.org/html/1909.02816}{Bischoff--Jones,
Lemma~3.1 and Corollary~3.4}, an algebra automorphism and its inverse
are convolution idempotents. They are positive sums of the primitive
idempotents $d_a[a]$; their fusion product is $[\mathbf1]$. Positivity
forces singleton supports consisting of inverse invertibles.

The earlier pointwise automorphism calculation gives
$\operatorname{Aut}(B/U)=1$. If rank six had six invertibles, then
$3s=6$, so $X$ would be invertible and $B=\mathbf1\oplus X$ would
have a nontrivial sign automorphism, a contradiction. Thus there are
exactly three invertibles in either rank. Their associator is also
determined: the local $B$-module category has metric group
\[
 (\mathbb Z/3)^2,\qquad q(a,b)=\frac{2a^2+b^2}{3}\pmod1.
\]
The actual extension $H/B$ selects an isotropic line. Forming its
modules gives an embedded untwisted $\operatorname{Vec}_{\mathbb Z/3}$
subcategory of $\mathcal F$. Merely knowing the group of invertibles
would not determine this associator.

\subsection{Excluding rank four: scope of the longer local argument}

At rank four the fusion ring is near-group:
\[
 a^3=1,\qquad a\rho=\rho a=\rho,\qquad
 \rho^2=1+a+a^2+m\rho.
\]
The published \href{https://doi.org/10.1016/j.jalgebra.2022.02.020}
{Dong--Chen--Wang, Theorem~6.2}, restricts $m$ to $0,2,3,6$.
For $m=0$, $s=2$ and the sign-automorphism contradiction just given
applies. For $m=2$, the categorical reduction of
\href{https://arxiv.org/html/math/0304194#S7}{Etingof--Gelaki--Ostrik,
Theorem~7.1 and Corollary~7.4} supplies the twisted $A_4$ double.
Its electric regular algebra gives a holomorphic $c=8$ extension
$H'$ with $U=(H')^{A_4}$, and
\href{https://msp.org/pjm/2004/213-2/pjm-v213-n2-p05-p.pdf}
{Dong--Mason, Theorem~1}, identifies $H'$ with $V_{E_8}$.
The order-two and order-three adjoint traces in $E_8$ are bounded
below by $-8$ and $-4$, respectively, as follows from the affine
Dynkin-diagram classification. Therefore
\[
 \dim(E_8^{A_4})\ge\frac{248+3(-8)+8(-4)}{12}=16,
\]
contradicting $\dim U_1=6$.

For clarity, the remaining exclusions use universal spherical
induction arguments, not a classification limited to unitary examples.
For $m=3k$, induction yields integers $p,q\ge0$ with $p+q=k$.
At $k=1$ the pointed factor and two indicator moments force the modular
representation of $G_{2,3}\otimes A_{2,1}$ or its reversal. Harmonic
theta-series pole cancellation in these two orientations gives
\[
 \dim U_1=15+A+10B\quad\hbox{or}\quad14+42A,
 \qquad A,B\in\mathbb Z_{\ge0}.
\]
At $k=2$, the balanced pattern $p=q=1$ forces, by signed Galois symmetry,
a row of the center's $S$ matrix whose Verlinde coefficient is
\[
 N_{XX}^{X}=\frac{97}{72}+\frac{32}{72}-\frac{33}{72}=\frac43.
\]
The unbalanced pattern $\{p,q\}=\{0,2\}$ admits a weighted character
projection. Its three-dimensional block is
\[
 S_v=-\frac12\begin{pmatrix}1&2&1\\1&0&-1\\1&-2&1\end{pmatrix},
 \qquad T_v=\operatorname{diag}(1,-i,-1)
\]
before the common central-charge factor; reversal conjugates $T_v$.
Harmonic theta-series subtraction gives the two alternatives
\[
 \dim U_1=8+A\quad\hbox{or}\quad5+12P,
 \qquad A,P\in\mathbb Z_{\ge0}.
\]
None of these formulas allows six currents. The parameters are actual
weighted polar top dimensions, so their integrality and positivity
do not require linear independence of all simple characters.

\begin{reviewnote}{R3: compressed rank-four exclusions}
The preceding paragraph is a summary, not a full proof of the universal
$m=3,6$ exclusions. The derivation of the complete modular representation
for $m=3$ is in
\artifact{hi_cft_20260923_next_m3_universal_data.md}, with an independent
audit in
\artifact{hi_cft_20260923_next_m3_universal_independent_audit.md}.
The $m=6$ arguments are in
\artifact{hi_cft_20260923_next_structure_m6_balanced_exclusion.md} and
\artifact{hi_cft_20260923_next_structure_m6_universal_projection.md}.
The coefficient calculations and pole-cancellation arguments are in
\artifact{hi_cft_20260923_next_structure_m3_character_obstruction.md}
and
\artifact{hi_cft_20260923_next_structure_m6_character_obstruction.md}.
These are new local deductions, not statements of the cited
rank-four classification theorem. Their universal induction reductions
and theta normalizations remain an explicit review dependency of this
condensed manuscript. In particular, the older character notes alone
only exclude specified model representations; the later universal
reductions are essential. The pointwise claim
$\operatorname{Aut}(B/U)=1$ and the $m=2$ electric-extension argument
also rely on the longer record
\artifact{hi_cft_20260923_next_structure_rank46.md} and its replacement
spherical audit
\artifact{hi_cft_20260923_next_structure_final_dependency_audit.md}.
\end{reviewnote}

\subsection{The rank-six ring and exclusion of higher multiplicity}

Rank four being excluded, \eqref{eq:rank-fork} gives
$H=\mathbf1\oplus X\oplus2Y$. The fusion algebra is
\[
 K_0(\mathcal F)\otimes\mathbb C
 \cong\mathbb C\oplus\mathbb C\oplus M_2(\mathbb C),
\]
by the composition algebra of $\operatorname{End}_{\mathcal C}(H)$.
Hence it is noncommutative. The three invertibles act transitively on
the three noninvertibles: otherwise both nontrivial idempotents of
$\mathbb C[\mathbb Z/3]$ split off, leaving a four-dimensional ideal
with a dimension character, incompatible with the $M_2$ block.
The left and right cyclic actions must be inverse; equality would
make the fusion ring commutative. Duality then fixes every noninvertible,
and reciprocity and associativity give
\begin{equation}
 a^3=1,\qquad\rho a=a^2\rho,\qquad
 \rho^2=1+m(\rho+a\rho+a^2\rho),\qquad m\ge1.
 \label{eq:rank6-ring}
\end{equation}

Here are the essential numerical identities in the local multiplicity
exclusion. Put
\[
 d=\frac{3m+\sqrt{9m^2+4}}2,\quad s=1+d^2,\quad
 \Delta=9m^2+4,\quad R=\rho+a\rho+a^2\rho.
\]
Center induction, with the untwisted pointed subcategory, yields
\[
 I(1)=1+X+2Y,\quad d_X=d^2,\quad d_Y=s,
 \qquad I(a)=I(a^2)=C_0+C_1+C_2,
\]
where $\theta_X=\theta_Y=1$, $d_{C_j}=s$, and the $C_j$ twists are
$1,\zeta_3,\zeta_3^2$.
All remaining central simples $Z_j$ have $F(Z_j)=n_jR$,
$n_j\in\mathbb Z_{>0}$. Global dimension and the first two
Ng--Schauenburg indicators give
\begin{align}
 \sum_jn_j^2&=4m^2+2,\notag\\
 \sum_jn_j^2\theta_j&=-\frac{m^2+m\sqrt\Delta}{2},
 \label{eq:indicator-first}\\
 \sum_jn_j^2\theta_j^2
 &=\frac{-m^2+(2\varepsilon-m)\sqrt\Delta}{2},
 \qquad\varepsilon=\nu_2(\rho)\in\{1,-1\}.
 \label{eq:indicator-second}
\end{align}
These follow from the center-twist trace formula in
\href{https://arxiv.org/pdf/math/0601012}{Ng--Schauenburg,
Theorem~4.1}. They concern a multiset of exactly $4m^2+2$ roots,
repeating $\theta_j$ precisely $n_j^2$ times.

Write $\Delta=t^2D$ with $D$ squarefree. One has $D\equiv1\pmod3$
and $7\nmid D$. Except when $D=10$, the elementary prime-factor estimate
$\varphi(2D)/\sqrt D\ge12/\sqrt{13}$ and
\href{https://arxiv.org/pdf/1401.1879}{Larson's root-count theorem}
applied to \eqref{eq:indicator-first} give
\[
 4m^2+2\ge\frac{mt}{2}\varphi(2D)
 \ge\frac{6m\sqrt{9m^2+4}}{\sqrt{13}}.
\]
For $m>1$ this is impossible, since
\[
 36m^2(9m^2+4)-13(4m^2+2)^2
 =4(m^2-1)(29m^2+13)>0.
\]
For $D=10$, a root multiset with sum $-a-b\sqrt{10}$, with
$a,b\in\mathbb Z_{\ge0}$, has at least $a+8b$ members. In the equality
case its squared-root sum is $a$, provided that sum lies in
$\mathbb Q(\sqrt{10})$. To see the equality condition, average over
the Galois group fixing $\sqrt{10}$. Rational contributions cost at
least one root per negative unit, with equality only for $-1$;
quadratic contributions cost at least eight roots per unit, with
equality only for a primitive-40 orbit. That orbit has squared sum
zero. The hypothesis on the second sum permits undoing the average.

In our application $m,t$ are even and $a=m^2/2$, $b=mt/2$.
The bound
\[
 4m^2+2\ge\frac{m^2}{2}+4mt
\]
excludes even $m\ge4$. The remaining $m=t=2$ is the equality case,
so the squared sum must be $2$, whereas \eqref{eq:indicator-second}
gives $-2$ or $-2-4\sqrt{10}$. Thus $m=1$.
The full conductor-orbit proof, including imprimitive orbits and the
second-moment hypothesis, is retained in
\artifact{hi_cft_20260923_next_structure_rank6_indicators.md} and
\artifact{hi_cft_20260923_next_rank6_root_audit.md}.

Consequently
\begin{equation}
 d=\frac{3+\sqrt{13}}2,\qquad
 s=\frac{13+3\sqrt{13}}2,\qquad
 d_X=\frac{11+3\sqrt{13}}2.
 \label{eq:haagerup-dimensions}
\end{equation}
For $m=1$, \eqref{eq:indicator-first} has six roots with sum
$(-1-\sqrt{13})/2$. The quadratic-period equality case forces the
six distinct nonresidue thirteenth roots, so $n_j=1$ and the center
has twelve simples. Its dimensions and twist exponents are
\[
 \begin{array}{c|cccc}
 \text{number}&1&1&4&6\\\hline
 \text{dimension}&1&d^2&s&3d
 \end{array}
 \qquad
 0,0,0,0,\tfrac13,\tfrac23,
 \tfrac2{13},\tfrac5{13},\tfrac6{13},
 \tfrac7{13},\tfrac8{13},\tfrac{11}{13}.
\]
\begin{reviewnote}{R4: conductor-orbit bounds and equality cases}
The reduction to $m=1$ and the count of twelve central simples use more
than the displayed indicator identities. The prime-factor estimate,
the exceptional $\mathbb Q(\sqrt{10})$ equality case, and the
quadratic-period equality case at $m=1$ are compressed here. Their full
conductor-orbit arguments, including imprimitive roots and multiplicities,
are in
\artifact{hi_cft_20260923_next_structure_rank6_indicators.md},
\artifact{hi_cft_20260923_next_rank6_root_audit.md}, and
\artifact{hi_cft_20260923_next_structure_rank6_center.md}.
These local deductions need review independently of the cited indicator
formula. The later net-sector count inherits this dependency.
\end{reviewnote}
The complete $S,T$ data follow from
\href{https://arxiv.org/html/2609.25185v1}{Gannon--Schopieray--Yadav,
Theorem~A(2), equivalently Theorem~5.1}: its relevant statement assumes
a complex spherical untwisted HI category. The auxiliary metric group
has prime order $13$ and its negative quadratic Gauss sum fixes its
isometry class. This identifies the standard Haagerup modular data.
It does not by itself identify a braided category; that stronger
conclusion is treated separately after complete unitarity.

\section{The complete character vector}
\label{sec:characters}

Let $q=e^{2\pi i\tau}$. Character exponents are $h-c/24=h-1/3$.
Using the preceding twist list and strict positivity of nonvacuum
lowest weights, the only allowed nonvacuum pole is
$\gamma q^{-7/39}$ in the twist-$2/13$ coordinate. Here
$\gamma\in\mathbb Z_{\ge0}$ is its top dimension if that module starts
at $2/13$, and is zero if its lowest weight is shifted upward. The
vacuum pole is $q^{-1/3}$ with coefficient one.

\begin{lemma}[Uniqueness from the principal part]
Two character vectors for this modular representation with the same
$\gamma$ coincide in every weight.
\end{lemma}
\begin{proof}
Their difference is holomorphic at infinity. The modular representation
has finite image, and covariance transports this property to every cusp
of a finite-index kernel. Each component is therefore a holomorphic
function on a compact modular curve, hence constant. A constant vector
must be fixed by both $S$ and the \emph{character} $T$ matrix, which
includes $e^{-2\pi i/3}$. The only $T$-fixed coordinate is the object
of categorical twist $e^{2\pi i/3}$. Its $S$ image has a nonzero vacuum
component, so it cannot span an invariant line. The constant is zero.
\end{proof}

\href{https://arxiv.org/pdf/1006.1326}{Evans--Gannon, Sections~1.4
and~2.4, equation~(2.7)}, provide a fundamental matrix $\Xi$ whose
vacuum and $d_3$ columns have precisely the two principal parts above.
Thus, in their labeling,
\[
 \boldsymbol\chi=\Xi_{:,0}+\gamma\Xi_{:,d_3},
 \qquad\dim U_1=6+13\gamma.
\]
The actual six-current result forces $\gamma=0$. Their recursive
construction of $\Xi$ and the uniqueness lemma give the entire vector,
not just a finite coefficient match. In particular,
\begin{equation}
 \operatorname{ch}U
 =q^{-1/3}\bigl(1+6q+120q^2+956q^3+6010q^4+\cdots\bigr).
 \label{eq:vacuum-character}
\end{equation}
The lowest weights and top dimensions are
\[
\begin{array}{c|c|r@{\qquad}c|c|r}
 \text{label}&h&\dim M_{\rm top}&\text{label}&h&\dim M_{\rm top}\\\hline
 0&0&1&d_1&6/13&7\\
 b&1&80&d_2&11/13&42\\
 a&1&81&d_3&15/13&119\\
 c_0&1&81&d_4&5/13&5\\
 c_1&1/3&3&d_5&7/13&13\\
 c_2&2/3&27&d_6&8/13&14
\end{array}
\]
In particular the final maximum is $15/13$, improving the preliminary
$17/10$ bound. The proof record is
\artifact{hi_cft_20260923_next_structure_conditional_characters.md}.
Although this vector fixes every unflavoured homogeneous dimension,
it does not determine the separate branching functions for the
fourteen-module core $A_L$, and it was not used to prove properness.
